\section{Processor and Measurement Methods}
\label{sec:methods}

\subsection{Latency and Throughput Definitions}
\label{sec:latency-throughput}

The execution time of a polynomial evaluator depends on two distinct
properties of the instruction stream.  The first is the delay introduced by
a dependency chain.  The second is the rate at which independent
instructions of the same class can be issued.

A \emph{dependency chain} is a sequence of operations in which each result
is required by the next operation.  The classical Horner recurrence is the
simplest example in this study:
\[
r_i=a_i+xr_{i+1}.
\]
Even if several multiplication units are available, the next multiplication
cannot begin before the preceding value \(r_{i+1}\) has been produced.

The \emph{latency} of an operation is the delay that it contributes to such
a dependency chain.  The \emph{throughput} is the rate at which independent
operations can be completed when enough independent operands are available.
For the algorithms considered here, neither quantity alone is sufficient.
A method with a longer local operation sequence can still execute faster if
many instances of that sequence are independent.

The Boolean-kernel reduction makes this distinction explicit.  At level
\(i\), all
\[
\frac{N}{2^{i+1}}
\]
folds are independent, while the transition from one level to the next is a
true dependency.  The critical path therefore grows with the number of
levels, \(m=\log_2 N\), rather than with the number of coefficients.

\subsection{Architecture Details}
\label{sec:architecture}

The measurements reported in this study were obtained on an Apple M1
processor.  The comparison follows the same general processor-oriented
principle used in earlier work on polynomial evaluation: arithmetic count
alone is insufficient when evaluation schemes expose different dependency
graphs \cite{ewart2020}.  The M1 is an out-of-order superscalar processor,
so several integer operations can be in flight simultaneously when the
instruction scheduler finds independent work.  Horner evaluation, blocked
Horner evaluation, and the Boolean-kernel tree expose very different
amounts of instruction-level and task-level parallelism.

We do not attempt to model the complete microarchitecture.  Only the
features required to interpret the measurements are considered: independent
integer multiplication, dependency depth, memory traffic, allocation, and
the scheduling of independent subtrees across CPU workers.

The field implementation also matters at the instruction level.  A
Goldilocks multiplication specialized to
\[
p=2^{64}-2^{32}+1
\]
uses one full-width product followed by shifts, additions, and
subtractions.  A Montgomery multiplication uses the initial product and an
additional product involving the reduction coefficient.  Figure
\ref{fig:field-dag} shows the corresponding dependency structure
schematically.

\begin{figure}[ht]
\centering
\setlength{\unitlength}{0.83mm}
\begin{picture}(150,65)

\put(5,59){\makebox(60,4)[c]{\small (a) Goldilocks special reduction}}
\put(95,59){\makebox(48,4)[c]{\small (b) Montgomery reduction}}

% Goldilocks DAG
\put(12,48){\framebox(18,7){\(a\times b\)}}
\put(12,35){\framebox(18,7){split hi/lo}}
\put(7,21){\framebox(28,7){split \(h_0,h_1\)}}
\put(7,8){\framebox(28,7){shift/add/sub}}
\put(12,48){\vector(0,-1){6}}
\put(21,35){\vector(0,-1){7}}
\put(21,21){\vector(0,-1){6}}

\put(43,48){\framebox(18,7){\(a+b\)}}
\put(43,35){\framebox(18,7){\(t(a+b)\)}}
\put(43,21){\framebox(18,7){reduce}}
\put(43,8){\framebox(18,7){\(+b\)}}
\put(52,48){\vector(0,-1){6}}
\put(52,35){\vector(0,-1){7}}
\put(52,21){\vector(0,-1){6}}

% Montgomery DAG
\put(94,48){\framebox(20,7){\(T=a b\)}}
\put(94,35){\framebox(20,7){\(m=T_0n'\)}}
\put(92,21){\framebox(24,7){\(T+mp\)}}
\put(94,8){\framebox(20,7){high word}}
\put(104,48){\vector(0,-1){6}}
\put(104,35){\vector(0,-1){7}}
\put(104,21){\vector(0,-1){6}}

\put(126,48){\framebox(20,7){\(\tilde a+\tilde b\)}}
\put(126,35){\framebox(20,7){MontMul}}
\put(126,21){\framebox(20,7){\(+\tilde b\)}}
\put(136,48){\vector(0,-1){6}}
\put(136,35){\vector(0,-1){7}}

\end{picture}
\caption{Schematic dependency graphs for the two field-arithmetic paths.
The left side shows the specialized Goldilocks product reduction and its
use in the kernel fold; the right side shows the Montgomery reduction path.
The diagram represents dependencies rather than exact machine
instructions.}
\label{fig:field-dag}
\end{figure}

\subsection{Exact Arithmetic}
\label{sec:exact-arithmetic}

Unlike floating-point polynomial evaluation, the present computation does
not involve a numerical approximation error.  Every benchmarked method is
required to return exactly the same element of \(\mathbb{F}_p\).

Correctness is checked before performance measurements.  The test suite
compares the optimized Goldilocks multiplication with a reference modular
multiplication, compares the fused Boolean-kernel fold with its direct field
expression, verifies the monomial-to-kernel change of basis by evaluating
the same polynomial in both representations, and verifies the Lagrange
implementation against the monomial reference.

The Montgomery implementation is subjected to the same requirement.
Canonical residues are converted to Montgomery form, multiplied there, and
converted back before comparison with the reference result.  This prevents
a performance result from being accepted merely because two implementations
use different internal encodings.

\subsection{Measurement Tools}
\label{sec:measurement-tools}

The benchmarks are implemented with Criterion.  The principal measurements
use a warm-up period before sampling and report an interval together with a
central estimate.  The benchmark harness applies \texttt{black\_box} to
inputs and results in order to prevent the compiler from replacing the
measured computation with a constant result.

Microbenchmarks are used for the field operations and the kernel fold.
Complete polynomial evaluators are measured separately.  This separation is
important: an improvement in the latency of one multiplication need not
translate linearly to the full evaluator once memory traffic, worker
scheduling, and allocation are included.

For cross-method comparisons, all implementations are compiled from the same
repository revision and executed on the same machine.  A ratio is formed
only from measurements collected under the same field backend and benchmark
configuration.

\subsection{Programming Model}
\label{sec:programming-model}

The implementation is written in Rust.  Parallel work is expressed with
Rayon.  The monomial baseline divides the coefficient vector into blocks,
evaluates those blocks independently, and combines the resulting values.
The Boolean-kernel implementation divides the reduction tree into
power-of-two subtrees.  Each worker reduces one complete subtree locally,
thereby avoiding a global synchronization after every tree level.

The optimized kernel evaluator additionally reuses worker-local scratch
storage.  The first tree level is written directly from the immutable input
slice into a half-size scratch buffer, after which the remaining local
levels are performed in place.

This implementation preserves the arithmetic graph derived in
Section~\ref{sec:kernel-fold}; it changes only the scheduling and memory
traffic of that graph.

\subsection{Processor and Compiler}
\label{sec:processor-compiler}

Table~\ref{tab:platform} records the platform used for the current
measurement campaign.  The values are generated from the machine on which
this manuscript is built.

\begin{table}[ht]
\centering
\caption{Processor, compiler, and benchmark environment.}
\label{tab:platform}
\begin{tabular}{ll}
\toprule
Item & Configuration\\
\midrule
Processor & Apple M1\\
Physical CPU cores & 8\\
Logical CPU cores & 8\\
Memory & 16 GiB\\
Rust compiler & rustc 1.96.0-nightly (55e86c996 2026-04-02)\\
Cargo & cargo 1.96.0-nightly (888f67534 2026-03-30)\\
Operating system & Darwin 25.6.0 arm64\\
Parallel runtime & Rayon\\
\bottomrule
\end{tabular}
\end{table}

The benchmarks are compiled in release mode.  The repository configuration
uses link-time optimization for the performance runs.  No result in the
paper is taken from a debug build.

\subsection{Validation}
\label{sec:validation}

A simple consistency check is available before examining the final results.
The classical Horner method forms one dependency chain, whereas the
Boolean-kernel evaluator exposes independent folds at every level.  The
latter should therefore benefit more strongly from multiple workers once
the problem size is large enough to amortize scheduling overhead.

At small sizes the opposite effect is possible: the overhead of creating
parallel tasks and managing scratch storage can be comparable to the useful
work.  The expected behavior is consequently a crossover rather than a
uniform speedup at all sizes.

The field backend provides a second validation point.  If the
specialized Goldilocks reduction has a lower isolated multiplication cost
than Montgomery reduction, this advantage should be visible most clearly in
the serial Horner benchmark, whose execution is dominated by a chain of
multiplications.  In the Boolean-kernel evaluator the effect can be reduced
or amplified by memory traffic and parallel scheduling.  The measurements
in the following section are interpreted against these two expectations.
